\documentclass[10pt,a4paper]{amsart}
\usepackage[margin=19mm]{geometry}
\usepackage{amsmath,amssymb,mathtools}
\usepackage[scr=rsfs,cal=euler]{mathalfa}
\usepackage{xfrac}
\usepackage[unicode,hidelinks,bookmarks=false]{hyperref}
\newcommand{\ie}{i.e.\ }
\newcommand{\cf}{cf.\ }
\newcommand{\resp}{respectively\ }
\newcommand{\Subsection}{\subsection}
\providecommand{\nobreakdash}{\nobreak\hbox{-}}
\setlength{\emergencystretch}{2em}
\multlinegap0pt
\usepackage{float}
\usepackage[shortlabels]{enumitem}
\usepackage{dsfont}
\usepackage{amssymb}
\usepackage{mathtools}
\usepackage{caption}
\newtheorem{prop}{Proposition}
\newcommand{\R}{\ensuremath{\mathbb{R}}}
\title{The 15 Puzzle and homological stability in the space direction}
\author{Jes\'us Gonz\'alez, Matthew Kahle,, Nicholas Wawrykow}
\date{Source: arXiv:2602.21300v1 (CC BY 4.0)}
\begin{document}
\maketitle

\noindent\emph{Source and licence.} The 15 Puzzle and homological stability in the space direction. Version arXiv:2602.21300v1 (CC BY 4.0), distributed by arXiv under the Creative Commons Attribution 4.0 International licence, http://creativecommons.org/licenses/by/4.0/, as recorded in the arXiv licence metadata for this exact version and shown on the abstract page https://arxiv.org/abs/2602.21300v1. Full licence notice: \texttt{licenses/arxiv-2602.21300-CC-BY-4.0.txt}.\par\smallskip

\noindent\textit{Editorial scope.} Counted unit: the article's prop space (source0.tex lines 797--807). The article's complete original proof of it is delivered with it (source0.tex lines 809--834, one \texttt{proof} environment of 26 lines). No mathematics is written, repaired or re-derived: the statement and the proof are the article's own lines, in the article's own notation, and no retained line is substituted or re-flowed.  Retained uncounted as the source-local context the proof needs: lines 649--655; lines 727--733; lines 781--783.  Nothing was omitted from any retained span, so the package carries no omission record.  No retained line contains a citation, so the wrapper carries no bibliography and no bibliography entry.  No retained line contains a figure environment, an image-inclusion command or drawing code, so no image is delivered and the package declares no asset dependency.  Editorial formatting only, in addition: an AMS \texttt{amsart} preamble that restates the article's own declarations and macros used by the retained lines, namely the environments \texttt{prop} on separate counters, together with the standard packages those lines need.  The retained environments print in this document as Proposition 1 in order of appearance, whereas the article prints the same items with the numbering of its own full document.  The complete native source of the article as distributed in the arXiv e-print for this exact version is bundled unchanged as \texttt{sources/195259/source0.tex}.
\begin{prop}\label{n le wh-2 is connected}
If $\min\{w, h\}\ge 2$ and $wh-n\ge 2$, the configuration space $SF_{n}(R_{w,h})$ is path-connected.
Therefore, the inclusion $SF_{n}(R_{w,h})$ into $F_{n}(\R^{2})$ induces an isomorphism
\[
H_{0}\big(SF_{n}(R_{w,h})\big)\cong H_{0}\big( F_{n}(\R^{2})\big).
\]
\end{prop}

\begin{prop}\label{connected single right configuration space}
If $\min\{w-1, h\}\ge 2, wh-n\ge 2$, and $0\le p+1\le h$, then $SF_{n}(R_{w,h}; i_{0}, \dots, i_{p})$ is path-connected.
In particular, the inclusion of $SF_{n}(R_{w,h}; i_{0}, \dots, i_{p})$ into $F_{n}(\R^{2}; i_{0}, \dots, i_{p})$ induces an isomorphism
\[
H_{0}\big(SF_{n}(R_{w,h}; i_{0}, \dots, i_{p})\big)\cong H_{0}\big(F_{n}(\R^{2}; i_{0}, \dots, i_{p})\big).
\]
\end{prop}

\begin{prop}\label{squares in the two rightmost columns}
There is a deformation retraction from $SF_{n}(R_{w,h};j_{0}, \dots, j_{s};i_{0}, \dots, i_{p})$ onto its subspace in which the squares $i_{0}, \dots, i_{p}$ are tangent to the right side of $R_{w,h}$, and the squares $j_{0}, \dots, j_{s}$ are in the two right-most columns of $R_{w,h}$.
\end{prop}

\begin{prop}\label{connected double right configuration space}
If $n$, $w$, $h$, $s$ and $p$ are such that $\min\{w-1,h\}\ge 3$, $wh-n\ge \max\{6, h-p-1\}$, 
%$\min\big\{wh-6, (w-1)h+p+1\big\}\ge n$
$0\le s\le 2$, and $0\le s+1+p+1\le h$, then $SF_{n}(R_{w,h}; j_{0}, \dots, j_{s};i_{0}, \dots, i_{p})$ is path-connected.
In particular, the inclusion of the square configuration space $SF_{n}(R_{w,h}; j_{0}, \dots, j_{s};i_{0}, \dots, i_{p})$ into the point configuration space $F_{n}(\R^{2}; j_{0}, \dots, j_{s};i_{0}, \dots, i_{p})$ induces isomorphisms
\begin{align*}
H_{0}\big(SF_{n}(R_{w,h}; j_{0}, \dots, j_{s};i_{0}, \dots, i_{p})\big)&\cong H_{0}\big(SF_{n}(R_{w,h}; j_{0}, \dots, j_{s}, i_{0}, \dots, i_{p})\big)\\
&\cong H_{0}\big(F_{n}(\R^{2}; j_{0},\dots, j_{s}, i_{0}, \dots, i_{p})\big)\\
&\cong H_{0}\big(F_{n}(\R^{2}; j_{0}, \dots, j_{s}; i_{0}, \dots, i_{p})\big).
\end{align*}
\end{prop}

\begin{proof}

We show that we can slide the squares in any configuration in $SF_{n}(R_{w,h}; j_{0}, \dots, j_{s};i_{0}, \dots, i_{p})$ to the target configuration in which the squares $j_{0}, \dots, j_{s}, i_{0},\dots, i_{p}$ are the top $s+1+p+1$ squares of the right-most column of $R_{w,h}$ in that order, and the free squares are arranged in lexicographical order from top to bottom, right to left, starting directly below $i_{p}$.

If $p=-1$, then we can apply Proposition \ref{connected single right configuration space} to prove this result.
Otherwise, by Proposition \ref{squares in the two rightmost columns}, we may assume that the squares $j_{0}, \dots, j_{s},i_{0}, \dots, i_{p}$ are in the two right-most columns of $R_{w,h}$.
If the squares $j_{0}, \dots, j_{s}$ are already in the right-most column, then proceed with the next step; otherwise, note that since $s+1+p+1\le h$, we may slide the squares $i_{0}, \dots, i_{p}$ up and down so that the squares immediately to the right the $j_{0}, \dots, j_{s}$ are unoccupied. 
As such, we may slide the squares $j_{0}, \dots, j_{s}$ into the right-most column of $R_{w,h}$.
Moreover, applying Proposition \ref{connected single right configuration space} to this configuration as an element of $SF_{n}(R_{w,h};\sigma)$, where $\sigma$ is the shuffle of $(j_{0},\dots, j_{s})$ into $(i_{0},\dots, i_{p})$, we can move these squares to be at the top of the right-most column of $R_{w,h}$.
Additionally, since $(w-1)h+p+1\ge n$, we may assume that the free squares are all in the $R_{w-1,h}$ in the left side of $R_{w,h}$.

Next, we consider two cases.
If $s+1+p+1=h$, then since $wh-n\ge 6$, Proposition \ref{n le wh-2 is connected} tells us that we may slide the $n-h\le(w-1)h-6$ free squares about the $R_{w-1,h}$ in the left side of $R_{w,h}$ such that the grid-squares immediately to the left of the $j_{0},\dots, j_{s}$ are unoccupied. 
Moreover, if there is an $i_{l}$ directly above one of the $j$, we can further arrange the free squares such that the grid-square immediately to its left is unoccupied.
As such, we may slide $j_{0}, \dots, j_{s}$ into the second right-most column of $R_{w,h}$ and up one square if there is an $i_{l}$ immediately above it, the offending $i_{l}$ down one square, and the $j_{0},\dots, j_{s}$ back to the right, moving the $j_{0},\dots, j_{s}$ above the offending $i_{l}$.
Repeating this process if necessary, we can arrange the $j_{0}, \dots, j_{s}$ to be at the top of the right-most column of $R_{w,h}$ and the $i_{0},\dots, i_{p}$ to be directly below them.
Applying Proposition \ref{n le wh-2 is connected} allows us to rearrange the free squares to get the target configuration.

If $s+1+p+1<h$, then since $(w-1)h+p+1\ge n$, there are at least $s+1$ empty squares in the $R_{w-1,h}$ in the left side of $R_{w,h}$. 
We can slide these free squares about so that the grid-square immediately to the left of each of the $j_{0},\dots, j_{s}$ is unoccupied.
Slide all the $i_{0},\dots, i_{p}$ that are below $j_{s}$ to the bottom of the right-most column of $R_{w,h}$, and let $i_{l}$ be the last of the $i_{0},\dots, i_{p}$ above any of the $j_{0}, \dots, j_{s}$.
Since $s+1+p+1<h$, there is at least one square gap between $j_{s}$ and $i_{l+1}$.
Slide the $j_{0}, \dots, j_{s}$ to the left, the $i_{l}$ down to the square directly above $i_{l+1}$, and then the $j_{0}, \dots, j_{s}$ back to the right.
Repeating this process if necessary, we can arrange the squares such that the $j_{0},\dots, j_{s}$ are above the $i_{0}, \dots, i_{p}$.
Applying Proposition \ref{n le wh-2 is connected}, allows us to rearrange the free squares to get the target configuration.
\end{proof}

\end{document}
